\section{Метод и теоретические основания}\label{sec:method}
\subsection{Оценивание полезности и решение при квоте}
Основная идея --- сначала оценить сравнительную полезность второго
ответа по доступным признакам, затем отделить обучение этой оценки
от выбора запросов при ограничении на число вызовов.
Используются фиксированные эксперты и полностью парные метки;
статистическая независимость их ошибок не предполагается.
Признаки рассчитываются до второго вызова. Обучающий и эксплуатационный
потоки совпадают либо механизм изменения состава обучения известен.

Для пакета из $N$ запросов с условными полезностями $g_i=g(z_i)$
и ровно $k$ разрешёнными вызовами задача имеет вид
\begin{equation}\label{eq:topk}
\max_{r_i\in\{0,1\},\,\sum_i r_i=k}\ \sum_{i=1}^{N}r_i g_i.
\end{equation}
Предполагается, что остальные запросы пакета не дают дополнительной
информации об исходе объекта сверх $z_i$; это выполняется, например,
при независимых запросах. Решение --- выбрать $k$ наибольших $g_i$:
если выбран меньший элемент вместо большего, их обмен не уменьшает
цель. Это стандартное свойство квотного ранжирования, а не новый
теоретический результат. При ограничении \emph{не более} $k$ вызовов
отрицательную полезность эскалировать не следует; наши эксперименты
используют именно равную квоту, что важно для интерпретации вреда.

Если для всех объектов $|\widehat g_i-g_i|\le\varepsilon$,
разность ожидаемой accuracy оптимального и оценённого top-$k$ не больше
$2k\varepsilon/N$. Действительно, добавление и вычитание сумм
$\widehat g_i$ по двум наборам даёт два слагаемых, каждое ограничено
$k\varepsilon$; промежуточная разность неположительна по оптимальности
набора для $\widehat g$. Это условная граница: малая средняя MSE сама
по себе не гарантирует такую равномерную ошибку или верное ранжирование
около границы квоты.

\subsection{Модели и критерии обучения}
\textbf{Direct} --- линейная ridge-регрессия метки $D$ по
стандартизованным признакам. Квадратичная потеря на естественном потоке
минимизируется условным средним $g(z)$ в неограниченном классе функций;
линейный класс даёт регуляризованное приближение.
\textbf{Events} --- две logistic-регрессии событий $D=1$ и $D=-1$;
их разность используется как score. Каждая голова минимизирует
бинарную log-loss с L2-регуляризацией. Головы не обязаны давать
согласованную трёхклассовую вероятность; это ограничение параметризации.

\textbf{Weak-error} оценивает ошибку первой модели и вычитает среднюю
ошибку второй на fit. Вычитание константы не меняет ранжирование,
но score не равен условной полезности при неоднородной ошибке второго
эксперта. \textbf{Confidence} использует единицу минус максимальная
вероятность первого классификатора; не обучается и не рассматривается
как откалиброванный прогноз $D$.
Прямая оценка и разность двух \emph{ridge}-прогнозов с одинаковыми
данными и регуляризацией тождественны по линейности оператора решения.
Поэтому такое дублирование не образует содержательного нового метода.

Пусть $u_i$ --- вес потери и $n$ --- число обучающих извлечений.
В direct минимизируется $n^{-1}\sum_i u_i(D_i-f_\theta(z_i))^2+
\lambda\|\theta_{\rm coef}\|^2$; в каждой logistic-голове ---
$n^{-1}\sum_i u_i\ell_i+\lambda\|\theta_{\rm coef}\|^2/2$.
Intercept не штрафуется. В коде это соответственно
$\alpha=n\lambda$ и $C=(n\lambda)^{-1}$.
Для ridge решается регуляризованная линейная система,
для logistic применяется L-BFGS; предупреждение о несходимости
останавливает новый эксперимент. Отсутствующее событие даёт нулевую
константу, единственный наблюдаемый класс --- соответствующую константу.
Общая сетка регуляризации настраивается по MSE на tuning-dev.

\subsection{Почему отбор исходов меняет обучающую цель}
Пусть $p_t(z)=P(D=t\mid z)$, $t\in\{-1,0,1\}$,
и обучающий поток $Q$ перевзвешивает исход множителем $a_t>0$.
Тогда по формуле Байеса
\begin{equation}\label{eq:selection_shift}
\begin{aligned}
q_t(z)&=\frac{a_t p_t(z)}{\sum_{u=-1}^{1}a_u p_u(z)},\\
g_Q(z)&=\frac{a_1p_1(z)-a_{-1}p_{-1}(z)}{\sum_{u=-1}^{1}a_u p_u(z)}.
\end{aligned}
\end{equation}
Невзвешенная регрессия теперь оценивает $g_Q$, а не $g_P$.
Усиление вредных исходов меняет и числитель, и нормировку,
поэтому может изменить знак и порядок полезности.
Даже одинаковое усиление двух ненейтральных исходов не обязательно
сохраняет ranking: при вероятностях $(p_1,p_{-1},p_0)=(0.2,0,0.8)$
и $(0.4,0.15,0.45)$ исходные полезности равны 0.20 и 0.25;
при множителях $(a_1,a_{-1},a_0)=(4,4,1)$ получаются
0.50 и приблизительно 0.377. Это иллюстрация известного механизма
смещения отбора \citep{nyberg2023exploring}, не заявление новой теоремы.

В контроле используется конечный размеченный пул размера $N_0$.
Каждая строка извлекается \emph{с возвращением} с известной
вероятностью $\rho_i=a_{D_i}/\sum_j a_{D_j}$.
Вес обратной вероятности извлечения (IPW) равен
$u_i=1/(N_0\rho_i)$. Для фиксированной функции потерь
\begin{equation}\label{eq:ipw}
\mathbb E_{I\sim\rho}[u_I\ell_I]
=\sum_{i=1}^{N_0}\rho_i\frac{\ell_i}{N_0\rho_i}
=\frac{1}{N_0}\sum_{i=1}^{N_0}\ell_i.
\end{equation}
Таким образом, восстанавливается риск исходного \emph{пула}.
Равенство не гарантирует несмещённости обученных коэффициентов,
оптимальности ранжирования или прироста accuracy на новом потоке.
Вес не нормируется по реализации; стандартизатор общий для
corrected/uncorrected и обучается на их одинаковых fit-строках.
Взвешивание может увеличивать вариативность потерь, а повторные извлечения
не создают новых уникальных примеров. Поэтому одновременно измеряются
число уникальных строк и $n_{\rm eff}=(\sum_i u_i)^2/\sum_i u_i^2$;
последнее характеризует неравномерность весов, а не число новых меток.

\subsection{Алгоритм обучения и выбора}
\noindent\textbf{Вход:} размеченные непересекающиеся fit-пул,
tuning-dev, selector-dev; фиксированные эксперты и признаки;
квота, состав извлечений и сетка регуляризации.
\begin{enumerate}
\item Рассчитать парные исходы в fit-пуле и вероятности извлечений.
Получить fit-строки, сохранить их идентификаторы и веса.
\item Оценить стандартизацию только по fit. Для каждого метода
и значения регуляризации обучить кандидата, при коррекции передавая
веса в функцию потерь.
\item Выбрать регуляризацию по минимальной MSE на tuning-dev;
при точном равенстве предпочесть более сильную регуляризацию.
\item На selector-dev оценить четыре natural-кандидата при той же квоте.
Выбрать максимальную accuracy; при точном равенстве предпочесть
confidence, затем direct, weak-error, events. Это фиксированное правило,
а не универсальная упорядоченность сложности методов.
\item Применить выбранный score к новому пакету и выбрать top-$k$
с общими независимыми ключами разрешения равенств. Не переобучать
роутер после выбора; test использовать только для оценки.
\end{enumerate}
\noindent\textbf{Выход:} стандартизатор, обученный score, параметр
регуляризации и выбранный метод. Для диагностического сравнения
сохраняются также все кандидаты.

Для $d$ признаков решение ridge имеет стоимость порядка
$O(nd^2+d^3)$, одна оценка функции/градиента logistic --- $O(nd)$;
полная стоимость зависит от числа итераций и проверок шага L-BFGS.
Сетка из четырёх значений умножает стоимость обучения на четыре.
Предсказание линейных голов стоит $O(Nd)$, сортировка пакета ---
$O(N\log N)$. Здесь $d=8$, поэтому повторное обучение роутеров
не требует повторного запуска дорогих экспертов.
Проверяемые следствия: сокращение train/test разрыва с ростом данных,
устранение смещения известной коррекцией и ненулевая ошибка выбора
метода на конечном selector-dev.
